\chapter{Deprecated Interfaces}
\mpitermtitleindex{deprecated interfaces}
\label{chap:deprecated}
 
\section{Deprecated since \texorpdfstring{\mpiiidoto/}{MPI-2.0}} 
\label{sec:deprecated:since20}

The following function is deprecated and is superseded by \mpifunc{MPI\_COMM\_CREATE\_KEYVAL} in \mpiiidoto/. 
The language independent definition 
of the deprecated function is the same as that of the new function,
  except for the function name and a different behavior in the
  C/Fortran language interoperability, see \sectionref{sec:misc-attr}.
The language bindings are modified.

\begin{funcdef}{MPI\_KEYVAL\_CREATE}{(\mbox{copy\_fn}, \mbox{delete\_fn}, \mbox{keyval}, \mbox{extra\_state})}
\funcarg{\IN}{copy\_fn}{Copy callback function for \mpiarg{keyval}}
\funcarg{\IN}{delete\_fn}{Delete callback function for \mpiarg{keyval}}
\funcarg{\OUT}{keyval}{key value for future access (integer)}
\funcarg{\IN}{extra\_state}{Extra state for callback functions}
\end{funcdef}
\mpicallbackindex{MPI\_Copy\_function}%
\mpicallbackindex{MPI\_Delete\_function}%
\mpibindmain{MPI\_Keyval\_create}{(\gb{}MPI\_Copy\_function~*copy\_fn, MPI\_Delete\_function~*delete\_fn, int~*keyval, void~*extra\_state)}
\mpifnewnonebind{For this routine, an interface within the \code{mpi\_f08} module was never defined.}
\mpicallbackindex{COPY\_FUNCTION}%
\mpicallbackindex{DELETE\_FUNCTION}%
\mpifbindmain{MPI\_KEYVAL\_CREATE}{(COPY\_FN, DELETE\_FN, KEYVAL, EXTRA\_STATE, IERROR) \fargs EXTERNAL \mbox{COPY\_FN}, \mbox{DELETE\_FN}\fnfinalarg{}INTEGER \mbox{KEYVAL}, \mbox{EXTRA\_STATE}, \mbox{IERROR}}


The \mpiarg{copy\_fn} function is invoked when a communicator is
duplicated by \mpifunc{MPI\_COMM\_DUP}.  \mpiarg{copy\_fn} should be
of type \mpicallbackskip{MPI\_Copy\_function}, which is defined as follows:

\mpicallbackmainindex{MPI\_Copy\_function}
%%HEADER
%%SKIP
%%ENDHEADER

\mpitypedefbindmain{MPI\_Copy\_function}{(\gb{}MPI\_Comm~oldcomm, int~keyval, void~*extra\_state, void~*attribute\_val\_in, void~*attribute\_val\_out, int~*flag)}

A Fortran declaration for such a function is as follows:

\mpifnewnonebind{For this routine, an interface within the \code{mpi\_f08} module was never defined.}

\mpicallbackmainindex{COPY\_FUNCTION}%
\mpifsubbindmain{COPY\_FUNCTION}{(OLDCOMM, KEYVAL, EXTRA\_STATE, ATTRIBUTE\_VAL\_IN, ATTRIBUTE\_VAL\_OUT, FLAG, IERR) \fargs INTEGER \mbox{OLDCOMM}, \mbox{KEYVAL}, \mbox{EXTRA\_STATE}, \mbox{ATTRIBUTE\_VAL\_IN}, \mbox{ATTRIBUTE\_VAL\_OUT}, \mbox{IERR}\fnfinalarg{}LOGICAL \mbox{FLAG}}

\mpiarg{copy\_fn} may be specified as
\mpiconstfuncmain{MPI\_NULL\_COPY\_FN} or
\mpiconstfuncmain{MPI\_DUP\_FN}
from either C or Fortran;
\mpiconstfuncskip{MPI\_NULL\_COPY\_FN}
is a function that does nothing other than return \mbox{\mpiarg{flag}\mpicode{ = 0}}
and \error{MPI\_SUCCESS}.
\mpiconstfuncskip{MPI\_DUP\_FN} is a simple-minded
copy function that sets \mbox{\mpiarg{flag}\mpicode{ = 1}},
returns the value of
\mpiarg{attribute\_val\_in} in \mpiarg{attribute\_val\_out}, and
returns \error{MPI\_SUCCESS}.
Note that \mpiconstfuncskip{MPI\_NULL\_COPY\_FN}
and \mpiconstfuncskip{MPI\_DUP\_FN} are also deprecated.

% The folowing declarations are needed for generating the deprecated
% list of functions in the annex.
% They are not used for printing here.
% Therefore, they are surrounded by a deleting macro.
\OnlyForAutomaticAnnexGeneration{%
\mpibindmain{MPI\_NULL\_COPY\_FN}{(\gb{}MPI\_Comm~oldcomm, int~keyval, void~*extra\_state, void~*attribute\_val\_in, void~*attribute\_val\_out, int~*flag)}
\mpifbindmain{MPI\_NULL\_COPY\_FN}{(OLDCOMM, KEYVAL, EXTRA\_STATE, ATTRIBUTE\_VAL\_IN, ATTRIBUTE\_VAL\_OUT, FLAG, IERR) \fargs INTEGER \mbox{OLDCOMM}, \mbox{KEYVAL}, \mbox{EXTRA\_STATE}, \mbox{ATTRIBUTE\_VAL\_IN}, \mbox{ATTRIBUTE\_VAL\_OUT}, \mbox{IERR}\fnfinalarg{}LOGICAL \mbox{FLAG}}
}%

\OnlyForAutomaticAnnexGeneration{%
\mpibindmain{MPI\_DUP\_FN}{(\gb{}MPI\_Comm~oldcomm, int~keyval, void~*extra\_state, void~*attribute\_val\_in, void~*attribute\_val\_out, int~*flag)}
\mpifbindmain{MPI\_DUP\_FN}{(OLDCOMM, KEYVAL, EXTRA\_STATE, ATTRIBUTE\_VAL\_IN, ATTRIBUTE\_VAL\_OUT, FLAG, IERR) \fargs INTEGER \mbox{OLDCOMM}, \mbox{KEYVAL}, \mbox{EXTRA\_STATE}, \mbox{ATTRIBUTE\_VAL\_IN}, \mbox{ATTRIBUTE\_VAL\_OUT}, \mbox{IERR}\fnfinalarg{}LOGICAL \mbox{FLAG}}
}%


%Note that the C version of these two functions assume that the
%callback functions follow the C prototype, while the corresponding
%Fortran version assumes the Fortran prototype.

Analogous to \mpiarg{copy\_fn} is a callback deletion function, defined
as follows.  The \mpiarg{delete\_fn} function is invoked when a communicator is
deleted by \mpifunc{MPI\_COMM\_FREE} or when a call is made explicitly
to \mpifunc{MPI\_ATTR\_DELETE}.  \mpiarg{delete\_fn} should be
of type \mpicallbackskip{MPI\_Delete\_function}, which is defined as follows:

\mpicallbackmainindex{MPI\_Delete\_function}
%%HEADER
%%SKIP
%%ENDHEADER


\mpitypedefbindmain{MPI\_Delete\_function}{(\gb{}MPI\_Comm~comm, int~keyval, void~*attribute\_val, void~*extra\_state)}

A Fortran declaration for such a function is as follows:

\mpifnewnonebind{For this routine, an interface within the \code{mpi\_f08} module was never defined.}

\mpicallbackmainindex{DELETE\_FUNCTION}%
\mpifsubbindmain{DELETE\_FUNCTION}{(COMM, KEYVAL, ATTRIBUTE\_VAL, EXTRA\_STATE, IERR) \fargs \fnfinalarg{}INTEGER \mbox{COMM}, \mbox{KEYVAL}, \mbox{ATTRIBUTE\_VAL}, \mbox{EXTRA\_STATE}, \mbox{IERR}}

% RAB: The predefined callbacks are "constant" "function" pointers, therefore in both indexes
\mpiarg{delete\_fn} may be specified as
\mpiconstfuncmain{MPI\_NULL\_DELETE\_FN}
from either C or Fortran;
\mpiconstfuncskip{MPI\_NULL\_DELETE\_FN} is a function that does nothing other
than return \error{MPI\_SUCCESS}.
Note that \mpiconstfuncskip{MPI\_NULL\_DELETE\_FN} is also deprecated.

% The folowing declarations are needed for generating the deprecated
% list of functions in the annex.
% They are not used for printing here.
% Therefore, they are surrounded by a deleting macro.
\OnlyForAutomaticAnnexGeneration{%
\mpibindmain{MPI\_NULL\_DELETE\_FN}{(\gb{}MPI\_Comm~comm, int~keyval, void~*attribute\_val, void~*extra\_state)}
\mpifbindmain{MPI\_NULL\_DELETE\_FN}{(COMM, KEYVAL, ATTRIBUTE\_VAL, EXTRA\_STATE, IERROR) \fargs \fnfinalarg{}INTEGER \mbox{COMM}, \mbox{KEYVAL}, \mbox{ATTRIBUTE\_VAL}, \mbox{EXTRA\_STATE}, \mbox{IERROR}}
}%


\deprecatedsep
The following function is deprecated and is superseded by \mpifunc{MPI\_COMM\_FREE\_KEYVAL} in \mpiiidoto/. 
The language independent definition 
of the deprecated function is the same as the new function,
except for the function name. 
The language bindings are modified.

\begin{funcdef}{MPI\_KEYVAL\_FREE}{(\mbox{keyval})}
\funcarg{\INOUT}{keyval}{Frees the integer key value (integer)}
\end{funcdef}
\mpibindmain{MPI\_Keyval\_free}{(\gb{}int~*keyval)}
\mpifnewnonebind{For this routine, an interface within the \code{mpi\_f08} module was never defined.}
\mpifbindmain{MPI\_KEYVAL\_FREE}{(KEYVAL, IERROR) \fargs \fnfinalarg{}INTEGER \mbox{KEYVAL}, \mbox{IERROR}}


\deprecatedsep
The following function is deprecated and is superseded by \mpifunc{MPI\_COMM\_SET\_ATTR} in \mpiiidoto/. 
The language independent definition 
of the deprecated function is the same as the new function,
except for the function name. 
The language bindings are modified.

\begin{funcdef}{MPI\_ATTR\_PUT}{(\mbox{comm}, \mbox{keyval}, \mbox{attribute\_val})}
\funcarg{\INOUT}{comm}{communicator to which attribute will be attached (handle)}
\funcarg{\IN}{keyval}{key value, as returned by \mpifunc{MPI\_KEYVAL\_CREATE} (integer)}
\funcarg{\IN}{attribute\_val}{attribute value}
\end{funcdef}
\mpibindmain{MPI\_Attr\_put}{(\gb{}MPI\_Comm~comm, int~keyval, void~*attribute\_val)}
\mpifnewnonebind{For this routine, an interface within the \code{mpi\_f08} module was never defined.}
\mpifbindmain{MPI\_ATTR\_PUT}{(COMM, KEYVAL, ATTRIBUTE\_VAL, IERROR) \fargs \fnfinalarg{}INTEGER \mbox{COMM}, \mbox{KEYVAL}, \mbox{ATTRIBUTE\_VAL}, \mbox{IERROR}}

\deprecatedsep
The following function is deprecated and is superseded by \mpifunc{MPI\_COMM\_GET\_ATTR} in \mpiiidoto/. 
The language independent definition 
of the deprecated function is the same as the new function,
except for the function name. 
The language bindings are modified.

\begin{funcdef}{MPI\_ATTR\_GET}{(\mbox{comm}, \mbox{keyval}, \mbox{attribute\_val}, \mbox{flag})}
\funcarg{\IN}{comm}{communicator to which attribute is attached (handle)}
\funcarg{\IN}{keyval}{key value (integer)}
\funcarg{\OUT}{attribute\_val}{attribute value, unless \mpiarg{flag} \mpicode{= false}}
\funcarg{\OUT}{flag}{\mpicode{true} if an attribute value was extracted; \mpicode{false} if no attribute is associated with the key}
\end{funcdef}
\mpibindmain{MPI\_Attr\_get}{(\gb{}MPI\_Comm~comm, int~keyval, void~*attribute\_val, int~*flag)}
\mpifnewnonebind{For this routine, an interface within the \code{mpi\_f08} module was never defined.}
\mpifbindmain{MPI\_ATTR\_GET}{(COMM, KEYVAL, ATTRIBUTE\_VAL, FLAG, IERROR) \fargs INTEGER \mbox{COMM}, \mbox{KEYVAL}, \mbox{ATTRIBUTE\_VAL}, \mbox{IERROR}\fnfinalarg{}LOGICAL \mbox{FLAG}}


\deprecatedsep
The following function is deprecated and is superseded by \mpifunc{MPI\_COMM\_DELETE\_ATTR} in \mpiiidoto/. 
The language independent definition 
of the deprecated function is the same as the new function,
except for the function name. 
The language bindings are modified.

\begin{funcdef}{MPI\_ATTR\_DELETE}{(\mbox{comm}, \mbox{keyval})}
\funcarg{\INOUT}{comm}{communicator to which attribute is attached (handle)}
\funcarg{\IN}{keyval}{The key value of the deleted attribute (integer)}
\end{funcdef}
\mpibindmain{MPI\_Attr\_delete}{(\gb{}MPI\_Comm~comm, int~keyval)}
\mpifnewnonebind{For this routine, an interface within the \code{mpi\_f08} module was never defined.}
\mpifbindmain{MPI\_ATTR\_DELETE}{(COMM, KEYVAL, IERROR) \fargs \fnfinalarg{}INTEGER \mbox{COMM}, \mbox{KEYVAL}, \mbox{IERROR}}


\section{Deprecated since \texorpdfstring{\mpiiidotii/}{MPI-2.2}} 
\label{sec:deprecated:since22}

\label{sec:deprecated-cxx-bindings}
The entire set of C++ language bindings was deprecated
as of \MPIIIDOTII/ and removed in \MPIIIIDOTO/.
See Chapter~\ref{chap:removed}, \nameref{chap:removed}
for more information.

\deprecatedsep
The following function typedefs have been deprecated and are 
superseded by new names.  Other than the typedef names, the function 
signatures are exactly the same; the names were updated to match 
conventions of other function typedef names. 
 
\begin{center} 
\begin{tabular}{ll} 
\textbf{Deprecated Name} & \textbf{New Name} \\ 
\hline 
\mpicallbackmain{MPI\_Comm\_errhandler\_fn} & \mpicallback{MPI\_Comm\_errhandler\_function} \\ 
\mpicallbackmain{MPI\_File\_errhandler\_fn} & \mpicallback{MPI\_File\_errhandler\_function} \\ 
\mpicallbackmain{MPI\_Win\_errhandler\_fn}  & \mpicallback{MPI\_Win\_errhandler\_function}  \\ 
\hline 
\end{tabular} 
\end{center}  


\section{Deprecated since \texorpdfstring{\mpiivdoto/}{MPI-4.0}}
\label{sec:deprecated:since40}

\label{sec:deprecated-send-cancel}
\mpitermindex{cancel}
Cancelling a send request by calling \mpifunc{MPI\_CANCEL} has been deprecated and may be removed in a future version of the \MPI/ specification.


\deprecatedsep
The following function is deprecated and is superseded by the new
\mpifunc{MPI\_INFO\_GET\_STRING} call in \mpiivdoto/. 

\cdeclindex{MPI\_Info}%
\begin{funcdef}{MPI\_INFO\_GET}{(\mbox{info}, \mbox{key}, \mbox{valuelen}, \mbox{value}, \mbox{flag})}
\funcarg{\IN}{info}{info object (handle)}
\funcarg{\IN}{key}{key (string)}
\funcarg{\IN}{valuelen}{length of value associated with \mpiarg{key} (integer)}
\funcarg{\OUT}{value}{value (string)}
\funcarg{\OUT}{flag}{\mpicode{true} if \mpiarg{key} defined, \mpicode{false} if not (logical)}
\end{funcdef}
\mpibindmain{MPI\_Info\_get}{(\gb{}MPI\_Info~info, const~char~*key, int~valuelen, char~*value, int~*flag)}
\mpifnewbindmain{MPI\_Info\_get}{(info, key, valuelen, value, flag, ierror) \fargs TYPE(MPI\_Info), INTENT(IN)\ ::\ \mbox{info}\fnarg{}CHARACTER(LEN=*), INTENT(IN)\ ::\ \mbox{key}\fnarg{}INTEGER, INTENT(IN)\ ::\ \mbox{valuelen}\fnarg{}CHARACTER(LEN=valuelen), INTENT(OUT)\ ::\ \mbox{value}\fnarg{}LOGICAL, INTENT(OUT)\ ::\ \mbox{flag}\fnfinalarg{}INTEGER, OPTIONAL, INTENT(OUT)\ ::\ \mbox{ierror}}
\mpifbindmain{MPI\_INFO\_GET}{(INFO, KEY, VALUELEN, VALUE, FLAG, IERROR) \fargs INTEGER \mbox{INFO}, \mbox{VALUELEN}, \mbox{IERROR}\fnarg{}CHARACTER*(*) \mbox{KEY}, \mbox{VALUE}\fnfinalarg{}LOGICAL \mbox{FLAG}}

This function retrieves the value associated with key in a previous call to 
\mpifunc{MPI\_INFO\_SET}. If such a key exists, it sets \mpiarg{flag} to \mpicode{true} 
and returns the value in \mpiarg{value},
otherwise it sets \mpiarg{flag} to \mpicode{false} and leaves
\mpiarg{value} unchanged. 
\mpiarg{valuelen} is the number of characters available in 
value. If it is less than the actual size of the value, the value is
truncated. In C, \mpiarg{valuelen} should
be one less than the amount of allocated space to allow
for the null terminator. 

If \mpiarg{key} is larger than \mpiconst{MPI\_MAX\_INFO\_KEY}, 
the call is erroneous. 

The function \mpifunc{MPI\_INFO\_GET} is allowed to be called at any time, following the description for \MPI/ functionality that is always available in \sectionref{subsec:dynamic:preinit_functions}.

\deprecatedsep
The following function is deprecated and is superseded by the new
\mpifunc{MPI\_INFO\_GET\_STRING} call in \mpiivdoto/. 

\begin{funcdef}{MPI\_INFO\_GET\_VALUELEN}{(\mbox{info}, \mbox{key}, \mbox{valuelen}, \mbox{flag})}
\funcarg{\IN}{info}{info object (handle)}
\funcarg{\IN}{key}{key (string)}
\funcarg{\OUT}{valuelen}{length of value associated with \mpiarg{key} (integer)}
\funcarg{\OUT}{flag}{\mpicode{true} if \mpiarg{key} defined, \mpicode{false} if not (logical)}
\end{funcdef}
\mpibindmain{MPI\_Info\_get\_valuelen}{(\gb{}MPI\_Info~info, const~char~*key, int~*valuelen, int~*flag)}
\mpifnewbindmain{MPI\_Info\_get\_valuelen}{(info, key, valuelen, flag, ierror) \fargs TYPE(MPI\_Info), INTENT(IN)\ ::\ \mbox{info}\fnarg{}CHARACTER(LEN=*), INTENT(IN)\ ::\ \mbox{key}\fnarg{}INTEGER, INTENT(OUT)\ ::\ \mbox{valuelen}\fnarg{}LOGICAL, INTENT(OUT)\ ::\ \mbox{flag}\fnfinalarg{}INTEGER, OPTIONAL, INTENT(OUT)\ ::\ \mbox{ierror}}
\mpifbindmain{MPI\_INFO\_GET\_VALUELEN}{(INFO, KEY, VALUELEN, FLAG, IERROR) \fargs INTEGER \mbox{INFO}, \mbox{VALUELEN}, \mbox{IERROR}\fnarg{}CHARACTER*(*) \mbox{KEY}\fnfinalarg{}LOGICAL \mbox{FLAG}}

Retrieves the length of the \mpiarg{value} associated with
\mpiarg{key}. If \mpiarg{key} is defined, \mpiarg{valuelen} is set
to the length of its associated value and \mpiarg{flag} is set to
\mpicode{true}.  If \mpiarg{key} is not defined, \mpiarg{valuelen} is not
touched and \mpiarg{flag} is set to \mpicode{false}. The length returned in 
C does not include the end-of-string character.

If \mpiarg{key} is larger than \mpiconst{MPI\_MAX\_INFO\_KEY}, 
the call is erroneous.

The function \mpifunc{MPI\_INFO\_GET\_VALUELEN} is allowed to be called at any time, following the description for \MPI/ functionality that is always available in \sectionref{subsec:dynamic:preinit_functions}.

\deprecatedsep
The following return code has been deprecated and is superseded by a new name in \MPIIVDOTO/.
\begin{center}
\begin{tabular}{ll}
\textbf{Deprecated Name} & \textbf{Replacement Name} \\
\hline
\error{MPI\_T\_ERR\_INVALID\_ITEM} & \error{MPI\_T\_ERR\_INVALID\_INDEX} \\
\hline
\end{tabular}
\end{center}


\deprecatedsep
\label{sec:deprecated-fortran-sizeof}
The following Fortran subroutines are deprecated because the Fortran
language \cfunc{storage\_size()} and \cfunc{c\_sizeof()} intrinsic
functions provide similar functionality.  Note that while
\mpifunc{MPI\_SIZEOF} and \cfunc{c\_sizeof()} return the size in
bytes, \cfunc{storage\_size()} provides the size in bits.
\mpitermtitleindex{MPI\_SIZEOF and storage\_size@\protect\mpifuncskip{MPI\_SIZEOF} and \protect\cfunc{storage\_size()}}

\begin{funcdef}{MPI\_SIZEOF}{(\mbox{x}, \mbox{size})}
\funcarg{\IN}{x}{a Fortran variable of numeric intrinsic type (choice)}
\funcarg{\OUT}{size}{size of machine representation of that type (integer)}
\end{funcdef}
\mpifnewbindmainref{MPI\_Sizeof}{(x, size, ierror) \fargs TYPE(*), DIMENSION(..)\ ::\ \mbox{x}\fnarg{}INTEGER, INTENT(OUT)\ ::\ \mbox{size}\fnfinalarg{}INTEGER, OPTIONAL, INTENT(OUT)\ ::\ \mbox{ierror}}
\mpifbindmain{MPI\_SIZEOF}{(X, SIZE, IERROR) \fargs <type> \mbox{X}\fnfinalarg{}INTEGER \mbox{SIZE}, \mbox{IERROR}}

This function returns the size in bytes of the machine representation of
the given variable. It is a generic Fortran routine and has a Fortran
binding only. 

\begin{users} 
This function is similar to the C \code{sizeof} operator
but behaves slightly differently. If given an array argument, it 
returns the size of the base element, not the size of the
whole array.
\end{users}  

\begin{rationale}
This function is not available in other languages because it would
not be useful. 
\end{rationale}

\section{Deprecated since \texorpdfstring{\mpiivdoti/}{MPI-4.1}}

\label{sec:deprecated:since41}

\label{sec:deprecated-mpif.h}
The use of the \code{mpif.h} include file has been deprecated.
Information supporting the transition to \code{USE mpi} or \code{USE mpi\_f08} is provided in Section~\ref{f90:basic}.

\deprecatedsep
\label{sec:deprecated-mpi-host}
The predefined attribute key \mpiconst{MPI\_HOST} for \mpiconst{MPI\_COMM\_WORLD} when using the \wpm{} is deprecated.

\mpiconstmain{MPI\_HOST}: Host process rank, if such exists, \mpiconst{MPI\_PROC\_NULL}, otherwise.

\subsubsection*{Host Rank}
\mpitermtitleindex{host rank}
The value returned for \mpiconst{MPI\_HOST} gets the rank of the \mpitermni{HOST} process in the group associated
with communicator \mpiconst{MPI\_COMM\_WORLD}, if there is such.
\mpiconst{MPI\_PROC\_NULL} is returned if there is no host.
\MPI/ does not specify what it
means for a process to be a \mpitermni{HOST}, nor does it requires that a \mpitermni{HOST}
exists.

The attribute \mpiconst{MPI\_HOST} has the same value on all processes
of \mpiconst{MPI\_COMM\_WORLD}.

\begin{center}
\begin{tabular}{l}
\multicolumn{1}{c}{\textbf{Environmental inquiry keys}} \\
\hline
{\small C type: \ctype{const int} (or unnamed \ctype{enum})} \\
{\small Fortran type: \ftype{INTEGER}} \\
\hline
\mpiconst{MPI\_HOST} \\
\hline
\end{tabular}
\end{center}

\deprecatedsep
All \mpiskipfunc{MPI\_\XXX/\_X} procedures have been deprecated and may be removed in a future version of the \MPI/ specification.
In the case of their C binding and their Fortran binding through the \code{mpi\_f08} module, they are superseded by the large count and large byte displacement bindings of their counterpart in the form of \mpiskipfunc{MPI\_\XXX/}.

\begin{funcdef}{MPI\_TYPE\_SIZE\_X}{(\mbox{datatype}, \mbox{size})}
\funcarg{\IN}{datatype}{datatype to get information on (handle)}
\funcarg{\OUT}{size}{datatype size (integer)}
\end{funcdef}
\mpibindmain{MPI\_Type\_size\_x}{(\gb{}MPI\_Datatype~datatype, MPI\_Count~*size)}
\mpifnewbindmain{MPI\_Type\_size\_x}{(datatype, size, ierror) \fargs TYPE(MPI\_Datatype), INTENT(IN)\ ::\ \mbox{datatype}\fnarg{}INTEGER(KIND=MPI\_COUNT\_KIND), INTENT(OUT)\ ::\ \mbox{size}\fnfinalarg{}INTEGER, OPTIONAL, INTENT(OUT)\ ::\ \mbox{ierror}}
\mpifbindmain{MPI\_TYPE\_SIZE\_X}{(DATATYPE, SIZE, IERROR) \fargs INTEGER \mbox{DATATYPE}, \mbox{IERROR}\fnfinalarg{}INTEGER(KIND=MPI\_COUNT\_KIND) \mbox{SIZE}}

The description of \mpifunc{MPI\_TYPE\_SIZE} is applicable to this deprecated \mpiskipfunc{MPI\_TYPE\_SIZE\_X} accordingly,
see \sectionref{subsec:pt2pt-addfunc}.

\cdeclindex{MPI\_Aint}%
\begin{funcdef}{MPI\_TYPE\_GET\_EXTENT\_X}{(\mbox{datatype}, \mbox{lb}, \mbox{extent})}
\funcarg{\IN}{datatype}{datatype to get information on (handle)}
\funcarg{\OUT}{lb}{lower bound of datatype (integer)}
\funcarg{\OUT}{extent}{extent of datatype (integer)}
\end{funcdef}
\mpibindmain{MPI\_Type\_get\_extent\_x}{(\gb{}MPI\_Datatype~datatype, MPI\_Count~*lb, MPI\_Count~*extent)}
\mpifnewbindmain{MPI\_Type\_get\_extent\_x}{(datatype, lb, extent, ierror) \fargs TYPE(MPI\_Datatype), INTENT(IN)\ ::\ \mbox{datatype}\fnarg{}INTEGER(KIND=MPI\_COUNT\_KIND), INTENT(OUT)\ ::\ \mbox{lb}, \mbox{extent}\fnfinalarg{}INTEGER, OPTIONAL, INTENT(OUT)\ ::\ \mbox{ierror}}
\mpifbindmain{MPI\_TYPE\_GET\_EXTENT\_X}{(DATATYPE, LB, EXTENT, IERROR) \fargs INTEGER \mbox{DATATYPE}, \mbox{IERROR}\fnfinalarg{}INTEGER(KIND=MPI\_COUNT\_KIND) \mbox{LB}, \mbox{EXTENT}}

The description of \mpifunc{MPI\_TYPE\_GET\_EXTENT} is applicable to this deprecated \mpiskipfunc{MPI\_TYPE\_GET\_EXTENT\_X} accordingly,
see \sectionref{subsec:pt2pt-extent}.

\cdeclindex{MPI\_Aint}%
\begin{funcdef}{MPI\_TYPE\_GET\_TRUE\_EXTENT\_X}{(\mbox{datatype}, \mbox{true\_lb}, \mbox{true\_extent})}
\funcarg{\IN}{datatype}{datatype to get information on (handle)}
\funcarg{\OUT}{true\_lb}{true lower bound of datatype (integer)}
\funcarg{\OUT}{true\_extent}{true extent of datatype (integer)}
\end{funcdef}
\mpibindmain{MPI\_Type\_get\_true\_extent\_x}{(\gb{}MPI\_Datatype~datatype, MPI\_Count~*true\_lb, MPI\_Count~*true\_extent)}
\mpifnewbindmain{MPI\_Type\_get\_true\_extent\_x}{(datatype, true\_lb, true\_extent, ierror) \fargs TYPE(MPI\_Datatype), INTENT(IN)\ ::\ \mbox{datatype}\fnarg{}INTEGER(KIND=MPI\_COUNT\_KIND), INTENT(OUT)\ ::\ \mbox{true\_lb}, \mbox{true\_extent}\fnfinalarg{}INTEGER, OPTIONAL, INTENT(OUT)\ ::\ \mbox{ierror}}
\mpifbindmain{MPI\_TYPE\_GET\_TRUE\_EXTENT\_X}{(DATATYPE, TRUE\_LB, TRUE\_EXTENT, IERROR) \fargs INTEGER \mbox{DATATYPE}, \mbox{IERROR}\fnfinalarg{}INTEGER(KIND=MPI\_COUNT\_KIND) \mbox{TRUE\_LB}, \mbox{TRUE\_EXTENT}}

The description of \mpifunc{MPI\_TYPE\_GET\_TRUE\_EXTENT} is applicable to this deprecated \mpiskipfunc{MPI\_TYPE\_GET\_TRUE\_EXTENT\_X} accordingly,
see \sectionref{subsec:pt2pt-true-extent}.

\label{func:mpi-get-elements-x}
\cdeclindex{MPI\_Status}%
\begin{funcdef}{MPI\_GET\_ELEMENTS\_X}{(\mbox{status}, \mbox{datatype}, \mbox{count})}
\funcarg{\IN}{status}{return status of receive operation (status)}
\funcarg{\IN}{datatype}{datatype used by receive operation (handle)}
\funcarg{\OUT}{count}{number of received basic elements (integer)}
\end{funcdef}
\mpibindmain{MPI\_Get\_elements\_x}{(\gb{}const~MPI\_Status~*status, MPI\_Datatype~datatype, MPI\_Count~*count)}
\mpifnewbindmain{MPI\_Get\_elements\_x}{(status, datatype, count, ierror) \fargs TYPE(MPI\_Status), INTENT(IN)\ ::\ \mbox{status}\fnarg{}TYPE(MPI\_Datatype), INTENT(IN)\ ::\ \mbox{datatype}\fnarg{}INTEGER(KIND=MPI\_COUNT\_KIND), INTENT(OUT)\ ::\ \mbox{count}\fnfinalarg{}INTEGER, OPTIONAL, INTENT(OUT)\ ::\ \mbox{ierror}}
\mpifbindmain{MPI\_GET\_ELEMENTS\_X}{(STATUS, DATATYPE, COUNT, IERROR) \fargs INTEGER \mbox{STATUS(MPI\_STATUS\_SIZE)}, \mbox{DATATYPE}, \mbox{IERROR}\fnfinalarg{}INTEGER(KIND=MPI\_COUNT\_KIND) \mbox{COUNT}}

The description of \mpifunc{MPI\_GET\_ELEMENTS} is applicable to this deprecated \mpiskipfunc{MPI\_GET\_ELEMENTS\_X} accordingly,
see \sectionref{subsec:pt2pt-datatypeuse}.


\label{func:mpi-status-set-elements-x}

\cdeclindex{MPI\_Status}%
\begin{funcdef}{MPI\_STATUS\_SET\_ELEMENTS\_X}{(\mbox{status}, \mbox{datatype}, \mbox{count})}
\funcarg{\INOUT}{status}{status with which to associate count (status)}
\funcarg{\IN}{datatype}{datatype associated with count (handle)}
\funcarg{\IN}{count}{number of elements to associate with status (integer)}
\end{funcdef}
\mpibindmain{MPI\_Status\_set\_elements\_x}{(\gb{}MPI\_Status~*status, MPI\_Datatype~datatype, MPI\_Count~count)}
\mpifnewbindmain{MPI\_Status\_set\_elements\_x}{(status, datatype, count, ierror) \fargs TYPE(MPI\_Status), INTENT(INOUT)\ ::\ \mbox{status}\fnarg{}TYPE(MPI\_Datatype), INTENT(IN)\ ::\ \mbox{datatype}\fnarg{}INTEGER(KIND=MPI\_COUNT\_KIND), INTENT(IN)\ ::\ \mbox{count}\fnfinalarg{}INTEGER, OPTIONAL, INTENT(OUT)\ ::\ \mbox{ierror}}
\mpifbindmain{MPI\_STATUS\_SET\_ELEMENTS\_X}{(STATUS, DATATYPE, COUNT, IERROR) \fargs INTEGER \mbox{STATUS(MPI\_STATUS\_SIZE)}, \mbox{DATATYPE}, \mbox{IERROR}\fnfinalarg{}INTEGER(KIND=MPI\_COUNT\_KIND) \mbox{COUNT}}

The description of \mpifunc{MPI\_STATUS\_SET\_ELEMENTS} is applicable to this deprecated \mpiskipfunc{MPI\_STATUS\_SET\_ELEMENTS\_X} accordingly,
see \sectionref{sec:ei-status}.


